\documentclass[12pt]{article}
\usepackage[T1]{fontenc}
\usepackage[utf8]{inputenc}
\usepackage{lmodern}
\usepackage{textcomp}
\usepackage{enumitem}
\usepackage{geometry}
\usepackage{graphicx}
\usepackage{amsmath, amssymb}
\geometry{margin=1in}
\begin{document}
\begin{center}
Philosophy - Graduate Level Assessment 6 \
Total Marks: 40 \
Answer all questions.
\end{center}

\section*{Multiple Choice (10 marks)}
\begin{enumerate}[label=\arabic*.]
\item " \_Ad crumenam\_ " is a specific kind of
\begin{enumerate}[label=(\alph*)]
    \item Slippery Slope
    \item Hasty Conclusion
    \item False sign
    \item False analogy
    \item Bandwagon Fallacy
\end{enumerate}
\item What does "Svetambara" mean?
\begin{enumerate}[label=(\alph*)]
    \item "Dark-clad"
    \item "Sun-clad"
    \item "Earth-clad"
    \item "Sky-clad"
    \item "Fire-clad"
\end{enumerate}
\item Nagel thinks that the core of the absolutist position is that
\begin{enumerate}[label=(\alph*)]
    \item all actions are morally equivalent.
    \item it is permissible to harm as a foreseen but unintended consequence of action.
    \item the ends justify the means.
    \item individuals should always act in their own best interest.
    \item the moral worth of an action is determined by its outcome.
\end{enumerate}
\item According to Lukianoff and Haidt, institutionalizing vindictive protectiveness will
\begin{enumerate}[label=(\alph*)]
    \item encourage students to think pathologically.
    \item ill-prepare them for the workforce.
    \item harm their ability to learn.
    \item all of the above.
\end{enumerate}
\item According to Reiman, van den Haag's argument leads to the conclusion that
\begin{enumerate}[label=(\alph*)]
    \item we should refrain from imposing the death penalty.
    \item the death penalty is a necessary evil.
    \item murder is wrong.
    \item we should institute death by torture.
    \item capital punishment serves as a deterrent.
\end{enumerate}
\end{enumerate}

\section*{True / False (10 marks)}
\textit{Answer True or False}
\begin{enumerate}[label=\arabic*.]
\setcounter{enumi}{5}
\item True or False: Philo considers the analogy used by Cleanthes to be strong and effective in supporting his argument.
\item True or False: An appeal to pride involves convincing someone of a claim by asserting that their intelligence makes them inherently right.
\item True or False: The statement demonstrates equivocation by using ambiguous language about the reservation status.
\item True or False: The fallacy of amphiboly is commonly referred to as reification, which involves treating abstract concepts as concrete entities.
\item True or False: Nagel argues that absolutism dictates we should never prevent murder under any circumstances.
\end{enumerate}

\section*{Long Form Response (20 marks)}
\textit{Answer each question with a clear and complete explanation.}
\begin{enumerate}[label=\arabic*.]
\setcounter{enumi}{10}
\item Analyze the given premises in propositional logic: \textasciitilde{}E ⊃ \textasciitilde{}F, G ⊃ F, H ∨ \textasciitilde{}E, H ⊃ I, \textasciitilde{}I. Discuss the immediate consequences of these premises and explain how you derive \textasciitilde{}H from them.
\item Discuss Augustine's assertion that evil cannot exist without love, analyzing the implications of this claim for understanding the nature of evil in philosophical thought.
\end{enumerate}

\vfill
\noindent\textit{End of Paper}
\end{document}